\documentclass[conference]{IEEEtran}

% -----------------------------------------------------------------------
% Required Packages
% -----------------------------------------------------------------------
\usepackage{graphicx}
\usepackage{amsmath}
\usepackage{amssymb}
\usepackage{listings}
\usepackage{xcolor}
\usepackage{hyperref}
\usepackage{cite}
\usepackage{booktabs}
\usepackage{algorithm}
\usepackage{algpseudocode}
\usepackage{url}
\usepackage{array}
\usepackage{multirow}
\usepackage{longtable}
\usepackage{tabularx}
\usepackage{enumitem}
\usepackage{float}

% -----------------------------------------------------------------------
% Listings Style Configuration
% -----------------------------------------------------------------------
\definecolor{codebg}{gray}{0.93}
\definecolor{codegreen}{rgb}{0,0.5,0}
\definecolor{codegray}{rgb}{0.4,0.4,0.4}
\definecolor{codeblue}{rgb}{0.0,0.0,0.7}

\lstdefinestyle{jsonStyle}{
  language=,
  backgroundcolor=\color{codebg},
  basicstyle=\ttfamily\small,
  keywordstyle=\color{codeblue}\bfseries,
  stringstyle=\color{codegreen},
  commentstyle=\color{codegray}\itshape,
  breaklines=true,
  frame=single,
  rulecolor=\color{gray!40},
  captionpos=b,
  morekeywords={mcpServers,command,args,env,type,name,servers,tools,inputSchema,properties,required},
}

\lstdefinestyle{pythonStyle}{
  language=Python,
  backgroundcolor=\color{codebg},
  basicstyle=\ttfamily\small,
  keywordstyle=\color{codeblue}\bfseries,
  stringstyle=\color{codegreen},
  commentstyle=\color{codegray}\itshape,
  breaklines=true,
  frame=single,
  rulecolor=\color{gray!40},
  captionpos=b,
}

\lstdefinestyle{tsStyle}{
  language=,
  backgroundcolor=\color{codebg},
  basicstyle=\ttfamily\small,
  keywordstyle=\color{codeblue}\bfseries,
  stringstyle=\color{codegreen},
  commentstyle=\color{codegray}\itshape,
  breaklines=true,
  frame=single,
  rulecolor=\color{gray!40},
  captionpos=b,
}

% -----------------------------------------------------------------------
% Hyperref Config
% -----------------------------------------------------------------------
\hypersetup{
  colorlinks=true,
  linkcolor=black,
  citecolor=black,
  urlcolor=blue
}

% -----------------------------------------------------------------------
% Document Begins
% -----------------------------------------------------------------------
\begin{document}

\title{Context Engineering via Model Context Protocol:\\
A Deep Technical Study of Agentic AI Orchestration in an AI Study Agent}

\author{
  \IEEEauthorblockN{Lakshya Sharma}
  \IEEEauthorblockA{
    School of Engineering and Technology\\
    Sharda University\\
    Greater Noida, India\\
    \texttt{lakshya4568@github.com}
  }
}

\maketitle

% -----------------------------------------------------------------------
% Abstract
% -----------------------------------------------------------------------
\begin{abstract}
Large language model systems are no longer single-prompt artifacts.
They are runtime systems.
They route requests.
They retrieve evidence.
They execute tools.
They carry memory.
And they fail when context is overgrown, stale, or misallocated.
This paper studies that exact transition through a detailed case analysis of an AI Study Agent that combines LangGraph orchestration, Retrieval-Augmented Generation, and Model Context Protocol tool integration.
The central argument is technical and direct.
Prompt engineering optimizes text.
Context engineering optimizes state, contracts, and token allocation over time.
Using the project documents as primary evidence, we reconstruct the architecture, define a context lifecycle model, formalize token-budget tradeoffs, and identify concrete breakpoints in routing, retrieval precision, and tool exposure.
We show why MCP improves modularity and server interoperability while still allowing major inefficiencies when tool schemas are injected indiscriminately or when interaction traces are not compressed.
We then derive design constraints for context-aware orchestration: slot-bounded context assembly, query-adaptive tool gating, retrieval compression, and explicit reason-act-observe traces.
The broader contribution is a systems lens.
MCP is not merely a tool transport.
When used correctly, it becomes a context control plane for agentic AI.
When used naively, it increases overhead while preserving opacity.
\end{abstract}

\begin{IEEEkeywords}
context engineering, model context protocol, retrieval-augmented generation, agentic AI, tool calling, long-context optimization, reasoning traces
\end{IEEEkeywords}

% -----------------------------------------------------------------------
% I. Introduction
% -----------------------------------------------------------------------
\section{Introduction}

Large language models can now answer, summarize, reason, and synthesize.
That statement is true.
It is also incomplete.
In practical systems, model quality is often not the bottleneck.
Context quality is.
The question is no longer only ``what should the prompt say?''
The question is ``what should enter the context window at this specific step, and what must stay out?''

This distinction matters more in agentic applications than in static chat.
An agent is expected to perform multi-step work.
It routes intent.
It retrieves evidence.
It decides whether to invoke tools.
It may write to memory.
It may call tools repeatedly.
Each of these actions introduces context fragments that can help or hurt the next step.
Some context fragments are essential.
Some are noisy.
Some are stale but still expensive.

The attached project documents establish this pressure clearly.
They describe a system where LangGraph controls routing and execution flow, where a Python RAG backend provides retrieval context, and where MCP introduces dynamic tool schemas and server discovery.
This is exactly the type of architecture where prompt engineering alone becomes insufficient.
The engineering problem shifts from phrase tuning to information governance.

The paper therefore treats context engineering as a primary discipline.
Not as a rhetorical rebrand.
As a systems requirement.
We define context engineering as the design of structures, policies, and transformations that govern what information the model receives at every inference step.
That includes retrieval selection, schema exposure, memory compression, and lifecycle pruning.
It is broader than prompt engineering.
It is operational.

The paper asks six research questions.

\textbf{RQ1:}
What counts as context in an agentic LLM system beyond prompt text?

\textbf{RQ2:}
How does MCP change the structure of context and not only tool invocation mechanics?

\textbf{RQ3:}
Where do token costs scale nonlinearly in long-running sessions?

\textbf{RQ4:}
Which failure modes arise when context freshness and route policy diverge?

\textbf{RQ5:}
How should MCP-enabled systems gate tool visibility under token constraints?

\textbf{RQ6:}
What architecture-level changes convert MCP from abstraction overhead into reasoning leverage?

The main claims are intentionally assertive.

First.
MCP improves modularity and interoperability, but does not guarantee reasoning quality unless context assembly policy is disciplined.

Second.
Unbounded history replay and broad schema exposure create quadratic session-level token costs in practice.

Third.
Routing that overweights the latest utterance and underweights state produces avoidable retrieval bypass and answer drift.

Fourth.
A slot-based context policy with adaptive schema gating offers a cleaner optimization target than prompt-only refinement.

The rest of the paper is structured as follows.
Section II situates the work in related literature and standards.
Section III defines the system model and context taxonomy.
Section IV presents a full architecture reconstruction from the provided material.
Section V develops the core context engineering framework via MCP.
Section VI introduces mathematical analyses for token, latency, and retrieval behavior.
Section VII provides a critical failure analysis.
Section VIII compares the design to ReAct, Toolformer, and Reflexion paradigms.
Section IX proposes concrete redesign patterns.
Section X discusses limitations.
Section XI outlines future directions.
Section XII closes with a systems-level conclusion.

% -----------------------------------------------------------------------
% II. Related Work and Framing
% -----------------------------------------------------------------------
\section{Related Work and Framing}
\label{sec:related}

\subsection{Prompt Engineering and Its Ceiling}

Prompt engineering has produced major gains in few-shot performance and reasoning style control.
Its strengths are real.
Instruction shaping can alter decomposition quality.
Demonstration selection can shift output format stability.
Reasoning prompts can reduce superficial errors.

Yet prompt engineering has a structural ceiling in agentic systems.
It optimizes text artifacts, not runtime context architecture.
A perfect instruction still degrades if paired with stale retrieval chunks.
A careful chain-of-thought style still wastes tokens if irrelevant tool schemas are always injected.
A strong system prompt still fails if route policy skips retrieval at the wrong step.

This paper does not reject prompt engineering.
It repositions it.
Prompt engineering remains a component.
Context engineering becomes the envelope around it.

\subsection{RAG in Multi-Step Agent Systems}

RAG reduces hallucination by grounding generation in retrieved evidence.
That logic is now standard.
What is less discussed in implementation papers is lifecycle mismatch between retrieval context and conversation context.
If retrieval snapshots are appended but never pruned, they become expensive historical residue.
If retrieval is triggered only by explicit keywords, follow-up references can miss grounding.
If chunking is fixed and retrieval count is fixed, context payload may be either too sparse or too bloated.

The provided research material emphasizes these exact tensions.
It also highlights token-aware chunking and hybrid retrieval as practical mitigation.
Still, retrieval quality alone is not enough.
How retrieval outputs are inserted, prioritized, and expired matters equally.

\subsection{MCP as Standardized Capability Layer}

MCP introduces a host-client-server protocol model with typed tool schemas and runtime capability discovery.
This is significant.
It allows one agent host to interoperate with heterogeneous tool backends through a shared contract model.
It also enables dynamic expansion of capabilities without changing core orchestration code.

But this same strength creates new context questions.
Should all discovered tools be made visible to the model all the time?
Should schemas be transformed into dense natural language instructions on every step?
Should tool contracts be summarized differently for planning versus invocation phases?

The protocol does not impose one answer.
Host policy does.
This is precisely where context engineering enters.

\subsection{Agentic Reasoning Paradigms as Comparative Anchors}

ReAct emphasizes explicit thought-action-observation cycles.
Toolformer emphasizes learned and selective tool calls.
Reflexion emphasizes iterative self-feedback memory.
These paradigms expose one shared insight.
Reasoning quality improves when intermediate states are explicit, auditable, and policy-governed.

The studied system partially aligns with these paradigms.
It has routing and tool invocation paths.
It has persistent memory support.
It has retrieval-grounded generation.
But it still centralizes too much work in monolithic query calls.
And it still allows context accumulation patterns that blur relevance boundaries over long sessions.

% -----------------------------------------------------------------------
% III. Problem Definition and Context Taxonomy
% -----------------------------------------------------------------------
\section{Problem Definition and Context Taxonomy}
\label{sec:taxonomy}

\subsection{Operational Definition of Context}

In this paper, context is defined as any token payload or structured artifact that can influence the next model inference.
That includes:

\begin{enumerate}[leftmargin=1.2em]
\item system policy messages,
\item user and assistant conversational history,
\item retrieved document passages,
\item tool schema definitions and examples,
\item tool outputs and execution traces,
\item memory summaries,
\item route metadata and control tags.
\end{enumerate}

A token is not context only because it exists.
It becomes context when the model can attend to it.
That distinction matters for engineering.
Stored state and injected state are not the same thing.

\subsection{Context Classes}

We classify context into five classes.

\textbf{C1: Policy Context.}
Role, style, constraints, safety directives.
Usually stable across turns.

\textbf{C2: Interaction Context.}
Recent messages, unresolved references, discourse state.
High volatility.

\textbf{C3: Evidence Context.}
Retrieved chunks, citations, metadata.
Query-dependent.

\textbf{C4: Capability Context.}
Tool schemas, parameter contracts, invocation affordances.
Environment-dependent.

\textbf{C5: Memory Context.}
Long-term user or session abstractions.
Slow-changing but often compressed.

The central engineering mistake in many systems is to flatten all five classes into a single append-only stream.
That flattening is convenient for implementation.
It is costly for relevance.

\subsection{Context Lifecycle Stages}

Each context class moves through a lifecycle.

\textbf{Stage 1: Creation.}
A context artifact is generated or retrieved.

\textbf{Stage 2: Transformation.}
The artifact is formatted, compressed, enriched, or normalized.

\textbf{Stage 3: Injection.}
The artifact is inserted into model-visible input.

\textbf{Stage 4: Retention.}
The artifact remains in memory or history for future steps.

\textbf{Stage 5: Expiration.}
The artifact is pruned, summarized, or excluded from future calls.

Most systems implement creation and injection well enough.
They under-implement expiration.
The documents provided for this study repeatedly point to this gap.

\subsection{Objective Function}

Let context payload at turn $t$ be represented as
\begin{equation}
\mathcal{C}_t = \{P_t, H_t, R_t, Z_t, M_t\},
\end{equation}
where $P$ is policy, $H$ is history, $R$ is retrieval evidence, $Z$ is tool schema payload, and $M$ is memory summary.

Let token size of each component be $\tau(\cdot)$.
Given budget $B_t$, the feasibility constraint is
\begin{equation}
\tau(P_t)+\tau(H_t)+\tau(R_t)+\tau(Z_t)+\tau(M_t) \le B_t.
\end{equation}

We define utility of context component $x$ as $U(x|q_t,s_t)$ with query $q_t$ and state $s_t$.
The optimization problem is
\begin{equation}
\max_{\mathcal{C}_t} \sum_{x\in\mathcal{C}_t} w_x U(x|q_t,s_t)
\quad \text{s.t. token budget and freshness constraints.}
\label{eq:objective}
\end{equation}

Prompt engineering mostly tunes $U(P_t|\cdot)$.
Context engineering optimizes the full Eq.~\ref{eq:objective}.

% -----------------------------------------------------------------------
% IV. Reconstructed System Architecture
% -----------------------------------------------------------------------
\section{Reconstructed System Architecture}
\label{sec:architecture}

\subsection{Top-Level Structure}

The studied system follows a layered structure:
frontend interaction,
agent orchestration,
MCP capability integration,
and RAG service support.

The user flow is straightforward at surface level.
A query arrives.
A route is selected.
Retrieval and tool paths may be triggered.
A response is generated.

Under the hood, however, context enters from multiple channels at different times.
This temporal mismatch is what drives many subtle failures.

\begin{figure}[H]
\centering
\fbox{\parbox{0.94\columnwidth}{
\centering
\vspace{0.2cm}
\textbf{PNG Placeholder: fig_architecture_overview.png}\\
Insert full three-layer architecture diagram here\\
(Frontend, LangGraph orchestration, MCP servers, RAG backend)\\
\vspace{0.2cm}
}}
\caption{System architecture placeholder for a high-level visual. Replace with fig\_architecture\_overview.png.}
\label{fig:arch_overview}
\end{figure}

\subsection{LangGraph-Oriented Control Flow}

The workflow described in the project analysis can be summarized as:
route,
optional retrieve,
query,
optional tool execution,
response.

At first glance this appears clean.
In practice, two points are critical.

Point one.
Route decisions that depend heavily on the latest utterance can mis-handle follow-up references.

Point two.
If query processing injects broad policy plus broad tool metadata while replaying long history, token cost rises faster than expected.

\begin{figure}[H]
\centering
\fbox{\parbox{0.94\columnwidth}{
\centering
\vspace{0.2cm}
\textbf{PNG Placeholder: fig_langgraph_flow.png}\\
Insert node-edge flow diagram here\\
(router, retrieve, query, tools, flashcard, memory)\\
\vspace{0.2cm}
}}
\caption{LangGraph control-flow placeholder. Replace with fig\_langgraph\_flow.png.}
\label{fig:langgraph_flow}
\end{figure}

\subsection{MCP Capability Layer}

The documents emphasize dynamic tool discovery and schema-driven invocation.
This is a major architectural advantage.
Tool servers can be added or removed with lower coupling.
Schema contracts become machine-readable and auditable.

However, dynamic discovery should not imply dynamic saturation.
Discovering 20 tools does not mean all 20 must be visible to the model for every query.
Selective exposure is part of context engineering, not protocol support.

\subsection{RAG Pipeline Layer}

The RAG service performs ingestion,
chunking,
embedding,
storage,
retrieval,
optional reranking,
and generation support.

The provided material notes token-based chunking and hybrid retrieval as deliberate choices.
These choices improve baseline retrieval behavior.
Still, the context engineering question remains:
how many retrieved chunks should be injected at each step,
and in what transformed format?

\begin{figure}[H]
\centering
\fbox{\parbox{0.94\columnwidth}{
\centering
\vspace{0.2cm}
\textbf{PNG Placeholder: fig_rag_pipeline.png}\\
Insert RAG ingestion and retrieval pipeline visual here\\
(loader, chunker, embedder, vector DB, retriever, reranker, generator)\\
\vspace{0.2cm}
}}
\caption{RAG pipeline placeholder. Replace with fig\_rag\_pipeline.png.}
\label{fig:rag_pipeline}
\end{figure}

% -----------------------------------------------------------------------
% V. Context Engineering via MCP
% -----------------------------------------------------------------------
\section{Context Engineering via MCP}
\label{sec:core}

\subsection{Why MCP Matters for Context, Not Just Tools}

MCP is often introduced as a standard for tool connectivity.
That is correct but incomplete.
In an agentic system, tool schemas are not neutral metadata.
They are context tokens.
They influence planning.
They influence call probability.
They influence output format choices.

When schemas are concise and selectively exposed, they improve execution precision.
When schemas are verbose and always-on, they can dilute attention and inflate latency.
Hence,
MCP becomes a context-control instrument.
Not just a transport protocol.

\subsection{Schema as Context Contract}

A tool schema defines:
name,
description,
parameter set,
required fields,
defaults,
constraints.

This structured shape provides two advantages over prose-only tool descriptions.

First,
validation can be externalized and deterministic.

Second,
the model can receive compact affordance information that is easier to parse than ad-hoc textual instructions.

But there is also risk.
If schema enrichment adds many generated examples and detailed parameter commentary for every tool at every step, the schema payload can dominate context budget.

\begin{lstlisting}[style=jsonStyle, caption={Example MCP tool contract structure with validation-relevant fields.}, label={lst:mcp_contract}]
{
  "name": "generate_quiz",
  "description": "Create quiz items from topic and level",
  "inputSchema": {
    "type": "object",
    "properties": {
      "topic": {
        "type": "string",
        "minLength": 3,
        "description": "Subject or chapter"
      },
      "difficulty": {
        "type": "string",
        "enum": ["easy", "medium", "hard"],
        "default": "medium"
      },
      "count": {
        "type": "integer",
        "minimum": 1,
        "maximum": 30,
        "default": 10
      }
    },
    "required": ["topic"]
  }
}
\end{lstlisting}

\subsection{Tool Gating Policy}

Given tool set $\mathcal{Z}$,
query $q_t$,
and state $s_t$,
we define adaptive tool exposure:
\begin{equation}
\mathcal{Z}_t = \{z \in \mathcal{Z} \mid g(z,q_t,s_t)=1\}.
\label{eq:tool_gating}
\end{equation}

The gate $g$ may be heuristic,
classifier-based,
or model-assisted.
The key property is boundedness.
$|\mathcal{Z}_t|$ must remain small unless evidence justifies expansion.

This policy directly reduces schema token load.
It also reduces accidental tool deliberation.

\subsection{Context Slot Design}

We define five slots.

\textbf{Slot A: Policy slot.}
Stable system behavior directives.

\textbf{Slot B: Interaction slot.}
Recent conversational turns and unresolved references.

\textbf{Slot C: Evidence slot.}
Retrieved passages and citations.

\textbf{Slot D: Capability slot.}
Gated tool schemas and invocation affordances.

\textbf{Slot E: Memory slot.}
Compressed long-term user and task context.

Each slot receives a token cap.
Caps are query-adaptive.
Unused budget can be reassigned.
This is basic resource scheduling.
That is exactly what context engineering should be.

\begin{figure}[H]
\centering
\fbox{\parbox{0.94\columnwidth}{
\centering
\vspace{0.2cm}
\textbf{PNG Placeholder: fig_context_slots.png}\\
Insert slot-based context budget diagram here\\
(Policy, Interaction, Evidence, Capability, Memory)\\
\vspace{0.2cm}
}}
\caption{Slot-based context assembly placeholder. Replace with fig\_context\_slots.png.}
\label{fig:context_slots}
\end{figure}

\subsection{Slot Budgeting Formulation}

Let slot token allocations be $b_P,b_H,b_R,b_Z,b_M$.
Then
\begin{equation}
b_P+b_H+b_R+b_Z+b_M \le B_t.
\end{equation}

Let relevance estimates be $\hat{u}_P,\hat{u}_H,\hat{u}_R,\hat{u}_Z,\hat{u}_M$.
A practical allocator can solve
\begin{equation}
\max_{b} \sum_{x\in\{P,H,R,Z,M\}} \hat{u}_x \log(1+b_x)
\quad \text{s.t. } b_x\ge0,\sum_x b_x\le B_t.
\label{eq:allocator}
\end{equation}

The logarithm regularizes monopolization.
No single slot consumes all budget unless utility estimates are extreme.

\subsection{Inference-Time Assembly Procedure}

\begin{algorithm}[H]
\caption{Adaptive Context Assembly with MCP-Aware Tool Gating}
\label{alg:adaptive_context}
\begin{algorithmic}[1]
\Require query $q_t$, state $s_t$, budget $B_t$, full tool set $\mathcal{Z}$
\State Estimate route intent and retrieval need from $(q_t,s_t)$
\State Retrieve evidence candidates and compute relevance scores
\State Compress evidence candidates into bounded evidence slot
\State Build interaction slot from recent turns plus unresolved references
\State Load memory summary only if personalization or continuity is needed
\State Gate tool set using Eq.~\ref{eq:tool_gating}
\State Allocate slot budgets using Eq.~\ref{eq:allocator}
\State Trim and normalize each slot to fit assigned budget
\State Assemble final model input in deterministic slot order
\State Invoke model and capture tool/trace outputs
\State Update state with compressed trace artifacts
\State \Return response and updated state
\end{algorithmic}
\end{algorithm}

% -----------------------------------------------------------------------
% VI. Mathematical Analysis
% -----------------------------------------------------------------------
\section{Mathematical Analysis of Cost and Quality Tradeoffs}
\label{sec:math}

\subsection{Retrieval Similarity and Fusion}

For query embedding $\mathbf{q}$ and chunk embedding $\mathbf{c}_i$,
\begin{equation}
s_i = \cos(\mathbf{q},\mathbf{c}_i)=\frac{\mathbf{q}\cdot\mathbf{c}_i}{\|\mathbf{q}\|\|\mathbf{c}_i\|}.
\end{equation}

Hybrid ranking in the studied material can be represented using reciprocal rank fusion:
\begin{equation}
S(d)=\frac{w_s}{k+r_s(d)}+\frac{w_k}{k+r_k(d)},
\label{eq:rrf_math}
\end{equation}
where $w_s$ and $w_k$ are semantic and keyword weights,
$k$ is an RRF stabilizer,
and $r_s,r_k$ are respective ranks.

If semantic weight dominates,
conceptual queries benefit.
But topical ambiguity can increase near-neighbor noise.
Hence reranking or compression remains necessary before evidence injection.

\subsection{Session-Level Token Growth}

Assume per-turn injected tokens:
policy $p$,
new interaction tokens $u_t$,
assistant output tokens $a_t$,
retrieval tokens $r_t$,
schema tokens $z_t$,
memory tokens $m_t$.

If the system replays full prior history each turn,
prompt tokens at turn $t$ are approximated by
\begin{equation}
T_t \approx p + z_t + m_t + \sum_{i=1}^{t-1}(u_i+a_i+\rho_i r_i),
\end{equation}
with $\rho_i\in\{0,1\}$ indicating whether retrieval artifacts persist in replay.

Cumulative session token traffic:
\begin{equation}
\mathrm{Cost}(n)=\sum_{t=1}^{n} T_t.
\end{equation}

If average per-turn contribution is bounded away from zero,
$\mathrm{Cost}(n)$ scales quadratically in $n$.
This is the hidden tax of replay-heavy design.

\subsection{Latency Decomposition}

Total latency can be decomposed as
\begin{equation}
L_t = L^{(t)}_{route} + L^{(t)}_{retrieve} + L^{(t)}_{rerank} + L^{(t)}_{assemble} + L^{(t)}_{llm} + L^{(t)}_{tool}.
\end{equation}

$L^{(t)}_{llm}$ is sensitive to token count.
Thus token optimization indirectly optimizes latency.
Broad schema exposure affects both quality and time.

\subsection{Error Propagation Model}

Let route accuracy be $p_r$,
retrieval precision be $p_e$,
and tool-call correctness be $p_z$.
Under a simplified independence assumption,
step success probability is
\begin{equation}
P_{succ} \approx p_r\cdot p_e\cdot p_z.
\end{equation}

This model is simplistic but useful.
If any factor degrades,
overall success falls quickly.
Route reliability is often underrated in practice.

\subsection{Context Freshness Metric}

Define freshness score for evidence chunk $e_i$ at turn $t$:
\begin{equation}
F(e_i,t)=\exp(-\lambda\Delta t_i)\cdot \mathbb{I}[\text{topic match}],
\end{equation}
where $\Delta t_i$ is age in turns.

A naive replay policy effectively sets weak decay,
allowing stale evidence to persist.
Freshness-aware pruning should enforce thresholding on $F$ before reinjection.

\begin{figure}[H]
\centering
\fbox{\parbox{0.94\columnwidth}{
\centering
\vspace{0.2cm}
\textbf{PNG Placeholder: fig_token_growth_curve.png}\\
Insert line chart of cumulative token cost vs turns\\
(compare naive replay vs slot-bounded policy)\\
\vspace{0.2cm}
}}
\caption{Token-growth visualization placeholder. Replace with fig\_token\_growth\_curve.png.}
\label{fig:token_growth}
\end{figure}

% -----------------------------------------------------------------------
% VII. Deep Critical Analysis
% -----------------------------------------------------------------------
\section{Deep Critical Analysis of the Studied Design}
\label{sec:critical}

\subsection{Observation A: Context Is Layered, But Governance Is Uneven}

The design already distinguishes retrieval,
tool schemas,
memory,
and dialogue state.
This is good architecture intent.

The problem is governance asymmetry.
Creation and injection are explicit.
Expiration and pruning are less explicit.
That asymmetry creates accumulation drift.

\subsection{Observation B: MCP Modularity Is Strong at Boundaries}

MCP clearly lowers coupling between capability providers and host orchestration.
This is a genuine engineering win.
Teams can ship servers independently.
The host can discover them dynamically.

But modularity at boundaries does not imply efficiency inside inference payload.
If every server contributes always-visible schemas,
inference context can degrade despite clean service boundaries.

\subsection{Observation C: Monolithic Query Calls Hide Intermediate Reasoning}

When tool interaction traces are collapsed into final response text,
post-hoc diagnosis becomes harder.
Was the error due to wrong tool choice?
Wrong argument shape?
Wrong observation interpretation?
Without explicit intermediate artifacts,
root-cause analysis is slower and less reliable.

\subsection{Observation D: Retrieval and Routing Coupling Needs State Awareness}

If route decisions overfit lexical cues in the last utterance,
follow-up turns can lose retrieval grounding.
This breaks down especially for anaphoric user language:
``that section'',
``the previous topic'',
``summarize it now''.

A state-aware router should consult recent retrieved topic context,
not only current text surface.

\subsection{Observation E: Memory Can Improve Continuity or Amplify Noise}

Memory integration is useful when compressed and scoped.
It is harmful when raw and over-broad.
Memory tokens are expensive.
Memory relevance is variable.
Memory injection requires strict intent gating.

\subsection{Observation F: Context Bloat Is Usually Silent Until It Is Expensive}

Systems often appear stable in short tests.
The cost curve emerges later.
As session depth increases,
latency creeps,
retrieval relevance falls,
and responses become less anchored.
This is a lifecycle problem.
Not a single-step bug.

% -----------------------------------------------------------------------
% VIII. Failure Modes and Break Conditions
% -----------------------------------------------------------------------
\section{Failure Modes and Break Conditions}
\label{sec:failures}

\subsection{Failure Mode 1: Instruction Dilution by System Stacking}

Break condition:
multiple system-level directives are replayed with limited pruning.

Observed symptom:
responses remain fluent but citation fidelity declines.
The model follows high-level style directives while underweighting local evidence constraints.

Why it happens:
attention is finite.
Directive competition is real.
Priority ambiguity emerges when too many policy-like fragments coexist.

\subsection{Failure Mode 2: Retrieval Bypass in Follow-Up Turns}

Break condition:
router relies primarily on latest utterance lexical form.

Observed symptom:
follow-up conceptual questions receive plausible but weakly grounded answers.

Why it happens:
anaphoric turns lack explicit retrieval trigger terms.
Without state-aware routing,
retrieval can be skipped when it is still required.

\subsection{Failure Mode 3: Tool Over-Deliberation}

Break condition:
large ungated tool schema payload is always visible.

Observed symptom:
extra tool speculation,
longer response time,
higher token use without commensurate utility.

Why it happens:
the model evaluates more affordances than needed.
Capability abundance becomes cognitive noise.

\subsection{Failure Mode 4: Trace Opacity in Multi-Step Tasks}

Break condition:
intermediate tool steps are not represented as explicit inspectable state transitions.

Observed symptom:
hard-to-debug wrong outputs in workflows requiring multiple calls and reflections.

Why it happens:
collapsed traces remove causal visibility.

\subsection{Failure Mode 5: Session Drift Under Prolonged Interaction}

Break condition:
append-only history without adaptive summarization.

Observed symptom:
rising latency,
inconsistent grounding,
topic bleed.

Why it happens:
stale tokens persist longer than their utility horizon.

\subsection{Failure Mode 6: Shared Retrieval Space Contention}

Break condition:
insufficient isolation in vector retrieval context across users or tasks.

Observed symptom:
cross-topic contamination,
irrelevant chunks in top-$k$ retrieval.

Why it happens:
index segmentation and metadata filtering are underutilized.

\begin{table}[H]
\caption{Failure Modes, Trigger Conditions, and Primary Mitigations}
\label{tab:failure_matrix}
\centering
\begin{tabular}{p{1.9cm}p{2.0cm}p{2.1cm}}
\toprule
Failure Type & Trigger Condition & Primary Mitigation \\
\midrule
Instruction dilution & System-message accumulation & Directive dedup + slot precedence \\
Retrieval bypass & Last-utterance routing bias & State-aware route gating \\
Tool over-deliberation & Always-on schema exposure & Query-adaptive MCP tool gating \\
Trace opacity & Collapsed intermediate calls & Explicit reason-act-observe logs \\
Session drift & Full history replay & Adaptive summarization + caps \\
Retrieval contamination & Weak corpus isolation & Metadata filters + index partitioning \\
\bottomrule
\end{tabular}
\end{table}

\begin{figure}[H]
\centering
\fbox{\parbox{0.94\columnwidth}{
\centering
\vspace{0.2cm}
\textbf{PNG Placeholder: fig_failure_heatmap.png}\\
Insert risk heatmap for failure modes vs session depth\\
\vspace{0.2cm}
}}
\caption{Failure-risk heatmap placeholder. Replace with fig\_failure\_heatmap.png.}
\label{fig:failure_heatmap}
\end{figure}

% -----------------------------------------------------------------------
% IX. Comparative Analysis Against Agent Paradigms
% -----------------------------------------------------------------------
\section{Comparative Analysis Against Agent Paradigms}
\label{sec:comparative}

\subsection{Compared with ReAct-Style Loops}

ReAct emphasizes explicit alternation:
reason,
act,
observe,
reason again.

The studied design includes route and tool structures,
but can still collapse execution into opaque response generation in parts of the flow.
Hence it captures some operational benefits,
while missing full trace transparency.

\subsection{Compared with Toolformer-Style Selectivity}

Toolformer framing highlights that tools should be called when expected utility exceeds cost.
The studied architecture has protocol support for tools,
but selective invocation policy remains under-specified in the source material.
Adaptive gating is therefore a key improvement path.

\subsection{Compared with Reflexion-Style Self-Feedback}

Reflexion-oriented systems maintain explicit self-improvement traces.
The studied system has memory support,
which is a useful substrate.
Yet a dedicated reflective update loop with controlled policy integration is not deeply articulated.
That leaves potential gains unrealized.

\subsection{Compared with Prompt-Chaining Baselines}

The architecture is richer than simple prompt chaining.
It has dynamic routing,
retrieval,
and protocolized tools.

Still,
without strict context budgeting,
this richer architecture can inherit a prompt-chain weakness:
incremental token bloat that silently lowers relevance density over time.

\begin{table}[H]
\caption{Paradigm-Level Comparison}
\label{tab:paradigm_comparison}
\centering
\begin{tabular}{p{2.1cm}p{1.0cm}p{1.0cm}p{1.0cm}}
\toprule
Paradigm & Trace Visibility & Tool Selectivity & Long-Session Stability \\
\midrule
Prompt chaining & Low & Low & Low \\
ReAct-style & High & Medium & Medium \\
Toolformer-style & Medium & High & Medium \\
Reflexion-style & Medium & Medium & High \\
Studied MCP design & Medium & Medium (potentially low if ungated) & Medium \\
\bottomrule
\end{tabular}
\end{table}

% -----------------------------------------------------------------------
% X. Design Improvements and Redesign Blueprint
% -----------------------------------------------------------------------
\section{Design Improvements and Redesign Blueprint}
\label{sec:improvements}

\subsection{Improvement 1: Deterministic Slot Scheduler}

Implement a scheduler that allocates token budget across policy,
interaction,
evidence,
capability,
and memory slots.

Rules:

\begin{enumerate}[leftmargin=1.2em]
\item Hard-cap each slot.
\item Reallocate unused budget dynamically.
\item Apply freshness decay before reinjection.
\item Deduplicate directives at injection time.
\end{enumerate}

Expected impact:
lower token waste,
more stable grounding,
better long-session behavior.

\subsection{Improvement 2: Two-Stage Tool Exposure}

Stage A:
lightweight intent screening to identify candidate tool clusters.

Stage B:
inject only candidate schemas plus minimal examples.

Expected impact:
smaller schema payload,
faster planning,
fewer accidental tool branches.

\subsection{Improvement 3: State-Aware Routing}

Route function should include:

\begin{enumerate}[leftmargin=1.2em]
\item latest utterance,
\item unresolved discourse references,
\item recent retrieval-topic summary,
\item last route outcome quality signal.
\end{enumerate}

Expected impact:
reduced retrieval bypass,
improved follow-up grounding.

\subsection{Improvement 4: Explicit Trace Artifacts}

Persist reason-act-observe traces in compact form.
Inject compressed trace snippets for multi-step continuation.
This preserves causal explainability with bounded cost.

\subsection{Improvement 5: Retrieval Compression Layer}

Before evidence injection:

\begin{enumerate}[leftmargin=1.2em]
\item rerank candidate chunks,
\item compress to citation-preserving statements,
\item enforce source diversity constraints.
\end{enumerate}

Expected impact:
higher relevance density per token.

\subsection{Improvement 6: Memory Injection Gate}

Memory should not be always injected.
Enable only when query intent or task continuity demands it.

Expected impact:
reduced irrelevant personalization noise,
lower token load.

\begin{figure}[H]
\centering
\fbox{\parbox{0.94\columnwidth}{
\centering
\vspace{0.2cm}
\textbf{PNG Placeholder: fig_redesign_blueprint.png}\\
Insert redesigned orchestration blueprint\\
(slot scheduler + route gate + tool gate + compression stage)\\
\vspace{0.2cm}
}}
\caption{Redesign blueprint placeholder. Replace with fig\_redesign\_blueprint.png.}
\label{fig:redesign_blueprint}
\end{figure}

\subsection{Improvement 7: Monitoring Metrics for Context Health}

Introduce runtime metrics:

\begin{enumerate}[leftmargin=1.2em]
\item context relevance ratio,
\item stale-token fraction,
\item schema payload share,
\item retrieval citation usage rate,
\item route correction frequency,
\item tool call utility score.
\end{enumerate}

These metrics make context engineering observable.
What is not measured is rarely controlled.

\subsection{Improvement 8: Evaluation Harness for Context Policies}

Build an offline replay harness with controlled scenarios.
Compare policy variants on:
quality,
latency,
token cost,
and failure incidence.

This turns context policy from intuition into evidence-driven engineering.

% -----------------------------------------------------------------------
% XI. Extended Discussion
% -----------------------------------------------------------------------
\section{Extended Discussion}
\label{sec:discussion}

\subsection{MCP Adoption Pattern in Real Projects}

Many teams first adopt MCP for practical interoperability.
That is sensible.
It reduces ad-hoc integration work.
It creates reusable capability endpoints.

The next phase should be context-aware governance.
Without it,
MCP remains beneficial but under-leveraged.
The protocol can carry much more architectural value when tied to selective exposure and state policy.

\subsection{Context Engineering as an Organizational Discipline}

The findings suggest context engineering should be owned explicitly.
Not split implicitly across prompting,
retrieval,
and tools teams.

A practical structure:

\begin{enumerate}[leftmargin=1.2em]
\item maintain a context contract catalog,
\item define slot budgets per product mode,
\item enforce freshness and expiration SLAs,
\item run periodic context-drift audits.
\end{enumerate}

This mirrors how mature teams handle API contracts.
Context is now an API surface for inference.

\subsection{Why This Matters Beyond Education Use Cases}

The case study domain is educational assistance.
The engineering pattern is general.
Any domain with retrieval,
tools,
and long sessions faces the same context dynamics.

Examples include enterprise copilots,
clinical evidence assistants,
research workflow agents,
and developer tooling agents.
The specific tools differ.
The context lifecycle problem does not.

\subsection{Security and Governance Considerations}

Context engineering intersects with security.
Over-broad tool exposure increases attack surface.
Over-broad memory injection increases leakage risk.
Stale retrieval artifacts can revive sensitive data unexpectedly.

Hence policy design should include:

\begin{enumerate}[leftmargin=1.2em]
\item capability least privilege,
\item intent-validated tool invocation,
\item memory redaction layers,
\item provenance tagging for retrieved evidence.
\end{enumerate}

\subsection{Tradeoff Summary}

There is no free design.
Compression can remove nuance.
Strict gating can miss helpful tools.
Aggressive pruning can reduce continuity.

The goal is not maximal reduction.
The goal is allocation quality.

% -----------------------------------------------------------------------
% XII. Limitations
% -----------------------------------------------------------------------
\section{Limitations}
\label{sec:limitations}

This paper is grounded in the attached project documentation and existing manuscript context.
It does not claim a new benchmark dataset.
It does not report large-scale live A/B tests.

The mathematical models are intentionally simplifying.
They clarify engineering trends.
They are not full probabilistic simulators.

The comparative discussion references major paradigms for conceptual anchoring,
not for strict controlled equivalence.

The recommendations are design-oriented.
They still require implementation-specific validation before production adoption.

These limitations are acceptable for the goal of this study:
a deep architectural analysis that reframes MCP-enabled systems through context engineering constraints.

% -----------------------------------------------------------------------
% XIII. Future Directions
% -----------------------------------------------------------------------
\section{Future Directions}
\label{sec:future}

\subsection{Direction A: Learned Slot Allocation}

Replace heuristic slot budgets with learned policies that optimize quality-cost objectives under varying query classes.

\subsection{Direction B: Adaptive Schema Distillation}

Automatically compress schema descriptions into query-specific affordance summaries.

\subsection{Direction C: Retrieval Freshness Controllers}

Use freshness-aware confidence scores to decide reinjection versus expiration.

\subsection{Direction D: Trace-Aware Self-Repair}

Leverage explicit reason-act-observe traces to trigger automatic correction policies after failed tool calls.

\subsection{Direction E: Multi-Agent MCP Topologies}

Explore coordinated specialist agents sharing MCP capability pools with explicit arbitration and budget negotiation.

\subsection{Direction F: Standardized Context Health Benchmarks}

Create reusable benchmark suites for context relevance,
staleness,
compression loss,
and schema overhead in MCP-enabled systems.

\begin{figure}[H]
\centering
\fbox{\parbox{0.94\columnwidth}{
\centering
\vspace{0.2cm}
\textbf{PNG Placeholder: fig_future_roadmap.png}\\
Insert research roadmap visual with milestones\\
\vspace{0.2cm}
}}
\caption{Future roadmap placeholder. Replace with fig\_future\_roadmap.png.}
\label{fig:future_roadmap}
\end{figure}

% -----------------------------------------------------------------------
% XIV. Conclusion
% -----------------------------------------------------------------------
\section{Conclusion}
\label{sec:conclusion}

This paper argued that context engineering,
not prompt engineering alone,
is the decisive systems problem for modern agentic AI.

Using the supplied project analyses,
we reconstructed an MCP-enabled RAG agent architecture and examined where context is created,
transformed,
injected,
and left to decay.

The conclusions are clear.

MCP is valuable.
It provides real modularity and protocolized capability integration.

But MCP is not automatically efficient.
Without selective schema exposure,
slot-bounded assembly,
and explicit trace governance,
it can carry overhead that weakens long-session reasoning quality.

The path forward is concrete.
Treat context as a scheduled resource.
Apply deterministic lifecycle policy.
Gate capabilities by intent.
Compress evidence before injection.
Persist only useful traces.
Measure context health continuously.

Agentic systems fail quietly when context policy is vague.
They improve rapidly when context policy is explicit.

That is the core finding.

% -----------------------------------------------------------------------
% XV. Appendix A: Extended Analytical Notes
% -----------------------------------------------------------------------
\section*{Appendix A: Extended Analytical Notes}

This appendix expands implementation-facing analysis points so the manuscript can directly support system redesign work.

\subsection*{A1. Context Audit Checklist}

\begin{enumerate}[leftmargin=1.2em]
\item Enumerate all model-visible context artifacts by source.
\item Compute token share of each artifact category over 50-turn sessions.
\item Label each artifact by utility horizon (one turn, short horizon, long horizon).
\item Record reinjection frequency and freshness decay profile.
\item Identify stale artifacts surviving beyond utility horizon.
\item Measure route decision error under anaphoric follow-up prompts.
\item Measure schema payload growth as tool count increases.
\item Measure response-time sensitivity to schema payload volume.
\item Track citation fidelity relative to retrieval payload size.
\item Track memory-injection utility by intent class.
\end{enumerate}

\subsection*{A2. Slot Policy Template}

\begin{lstlisting}[style=pythonStyle, caption={Pseudo-template for context slot policy implementation.}, label={lst:slot_template}]
def allocate_slots(query, state, budget):
    profile = classify_query_profile(query, state)

    caps = {
        "policy": 0.10 * budget,
        "history": 0.25 * budget,
        "evidence": 0.35 * budget,
        "capability": 0.20 * budget,
        "memory": 0.10 * budget,
    }

    caps = adapt_caps(caps, profile)

    slots = {}
    slots["policy"] = build_policy_slot(state, caps["policy"])
    slots["history"] = build_history_slot(state, caps["history"])
    slots["evidence"] = build_evidence_slot(query, state, caps["evidence"])
    slots["capability"] = build_capability_slot(query, state, caps["capability"])
    slots["memory"] = build_memory_slot(query, state, caps["memory"])

    slots = normalize_and_trim(slots, budget)
    return slots
\end{lstlisting}

\subsection*{A3. Route Reliability Stress Cases}

Case 1:
Direct explicit retrieval request.
Expected route is retrieval-enabled.

Case 2:
Anaphoric follow-up with no explicit retrieval keyword.
Expected route should still preserve retrieval.

Case 3:
Tool request with mixed conceptual phrasing.
Expected route should gate to capability slot while preserving evidence continuity.

Case 4:
Cross-topic pivot mid-session.
Expected route should trigger stale evidence pruning and topic reset handling.

\subsection*{A4. Schema Payload Compression Heuristics}

\begin{enumerate}[leftmargin=1.2em]
\item Include full schema only for selected tools.
\item Convert non-selected tools to one-line affordance labels.
\item Limit examples to one per selected tool unless ambiguity detected.
\item Strip verbose prose from unchanged schema descriptions.
\item Cache and reference schema signatures by hash in trace logs.
\end{enumerate}

\subsection*{A5. Evidence Compression Heuristics}

\begin{enumerate}[leftmargin=1.2em]
\item Keep source identifiers and citation anchors intact.
\item Preserve at least one exact quote per high-impact source.
\item Remove duplicate semantic statements across chunks.
\item Penalize overly similar chunks from same source region.
\item Enforce minimum cross-source diversity when available.
\end{enumerate}

\subsection*{A6. Suggested Evaluation Matrix}

\begin{longtable}{p{2.0cm}p{2.4cm}p{2.6cm}}
\toprule
Dimension & Metric & Target Direction \\
\midrule
Token efficiency & Avg prompt tokens per turn & Decrease \\
Latency & P50 end-to-end latency & Decrease \\
Grounding & Citation-faithful answer ratio & Increase \\
Tool quality & Useful tool-call ratio & Increase \\
Route quality & Correct route selection rate & Increase \\
Stability & Drift incidents per 100 turns & Decrease \\
Memory utility & Helpful memory injection rate & Increase \\
Schema overhead & Schema token share & Decrease \\
Evidence quality & Relevant evidence density & Increase \\
Debuggability & Trace completeness score & Increase \\
\bottomrule
\end{longtable}

% -----------------------------------------------------------------------
% XVI. Appendix B: Placeholder Figure Index
% -----------------------------------------------------------------------
\section*{Appendix B: PNG Placeholder Figure Index}

For convenience,
all figure placeholders included in this manuscript are listed below.

\begin{enumerate}[leftmargin=1.2em]
\item fig\_architecture\_overview.png
\item fig\_langgraph\_flow.png
\item fig\_rag\_pipeline.png
\item fig\_context\_slots.png
\item fig\_token\_growth\_curve.png
\item fig\_failure\_heatmap.png
\item fig\_redesign\_blueprint.png
\item fig\_future\_roadmap.png
\end{enumerate}

These placeholders are currently rendered as framed text boxes to keep the source compilable before assets are prepared.

% -----------------------------------------------------------------------
% XVII. References
% -----------------------------------------------------------------------
\begin{thebibliography}{99}

\bibitem{anthropic_mcp}
Anthropic,
``Introducing the Model Context Protocol,''
2024.
[Online].
Available:
\url{https://www.anthropic.com/news/model-context-protocol}

\bibitem{mcp_spec}
Model Context Protocol Working Group,
``Model Context Protocol Specification,''
2025.
[Online].
Available:
\url{https://modelcontextprotocol.io/specification/2025-06-18}

\bibitem{anthropic_context}
Anthropic,
``Effective Context Engineering for AI Agents,''
2025.
[Online].
Available:
\url{https://www.anthropic.com/engineering/effective-context-engineering-for-ai-agents}

\bibitem{langchain_context}
LangChain,
``Context Engineering Documentation,''
2025.
[Online].
Available:
\url{https://docs.langchain.com/oss/python/langchain/context-engineering}

\bibitem{react}
S. Yao,
J. Zhao,
D. Yu,
N. Du,
I. Shafran,
K. Narasimhan,
and Y. Cao,
``ReAct: Synergizing Reasoning and Acting in Language Models,''
arXiv:2210.03629,
2022.

\bibitem{toolformer}
T. Schick,
J. Dwivedi-Yu,
R. Dessi,
R. Raileanu,
M. Lomeli,
E. Hambro,
L. Zettlemoyer,
N. Cancedda,
and T. Scialom,
``Toolformer: Language Models Can Teach Themselves to Use Tools,''
arXiv:2302.04761,
2023.

\bibitem{reflexion}
N. Shinn,
B. Labash,
A. Gopinath,
K. Narasimhan,
and S. Yao,
``Reflexion: Language Agents with Verbal Reinforcement Learning,''
arXiv:2303.11366,
2023.

\bibitem{context_doc}
L. Sharma,
``Context Engineering Notes (project document),''
\emph{AI Study Agent repository},
2026.

\bibitem{analysis_doc}
L. Sharma,
``Research Analysis Notes (project document),''
\emph{AI Study Agent repository},
2026.

\bibitem{j_split}
A. Arora et al.,
``JSPLIT: Prompt Bloating Mitigation for MCP Tool Catalogs,''
arXiv:2510.14537,
2025.

\bibitem{readme_ai}
X. Liu et al.,
``Readme\_AI: Context Servers for Agentic Workflows,''
IEEE Access,
2025.

\bibitem{rag_foundation}
P. Lewis et al.,
``Retrieval-Augmented Generation for Knowledge-Intensive NLP Tasks,''
NeurIPS,
2020.

\bibitem{dpr}
V. Karpukhin et al.,
``Dense Passage Retrieval for Open-Domain Question Answering,''
EMNLP,
2020.

\bibitem{fewshot}
T. B. Brown et al.,
``Language Models are Few-Shot Learners,''
NeurIPS,
2020.

\bibitem{cot}
J. Wei et al.,
``Chain-of-Thought Prompting Elicits Reasoning in Large Language Models,''
NeurIPS,
2022.

\bibitem{gpt4}
OpenAI,
``GPT-4 Technical Report,''
arXiv:2303.08774,
2023.

\bibitem{weaviate_context}
Weaviate,
``Context Engineering for Production LLM Systems,''
2025.

\bibitem{mcp_intro_blog}
P. Schmid,
``An Introduction to MCP for LLM Engineers,''
2025.

\bibitem{context_window_limits}
S. Orange,
``Prompt Length vs Context Window: Practical Limits in LLM Systems,''
2025.

\bibitem{mcp_github}
Model Context Protocol,
``Official MCP Repository,''
GitHub,
2025.
[Online].
Available:
\url{https://github.com/modelcontextprotocol}

\end{thebibliography}

\end{document}
